\begin{helper}[On the inner region, the derivative of a mollification is the mollified weak gradient]
Let $n\in\N$, let $\Omega\subseteq\R^n$, let $\varphi\colon\R^n\to\R$, let
$g_\varphi\colon\R^n\to\R^n$, let $\rho\colon\R^n\to\R$, let $d\in\R$, let
$x\in\R^n$ and let $i$ be an index in $\{1,\dots,n\}$.  Throughout, $|\cdot|$ is
the Euclidean norm, $D_i$ is the derivative in the direction of the $i$-th
standard basis vector, $\operatorname{supp}\rho$ denotes the \emph{closed}
support of $\rho$,
\[
(\varphi*\rho)(x)=\int_{\R^n}\varphi(y)\rho(x-y)\,dy
\]
is the convolution, and $\mathbf1_\Omega h$ denotes the zero extension of $h$,
i.e.\ the function equal to $h$ on $\Omega$ and to $0$ off $\Omega$.  Assume
\begin{enumerate}
\item[(H1)] $\Omega$ is open;
\item[(H2)] $g_\varphi$ is a weak gradient of $\varphi$ on $\Omega$
  (\noderef{HasWeakGradientOn}), and $\varphi$ is locally integrable on all of
  $\R^n$;
\item[(H3)] $\rho$ is $C^\infty$ with compact support and
  $\operatorname{supp}\rho=\overline{B_d(0)}$;
\item[(H4)] $\overline{B_d(x)}\subseteq\Omega$.
\end{enumerate}
Then
\[
D_i(\varphi*\rho)(x)\ =\ \bigl(\mathbf1_\Omega(D_i\varphi)\bigr)*\rho\,(x),
\]
where the right-hand side means the convolution of
$y\mapsto\mathbf1_\Omega(y)\,g_\varphi(y)_i$ against $\rho$.

The zero extension on the right is not cosmetic: (H2) controls $g_\varphi$ only
through an identity against test functions supported in $\Omega$ and says
nothing about $\varphi$ or $g_\varphi$ off $\Omega$, so the extension by zero is
what makes the right-hand convolution a well-defined global object.  The local
integrability of $\varphi$ in (H2), by contrast, is a hypothesis on all of
$\R^n$ and is what makes the left-hand convolution differentiable at all.
\end{helper}
\begin{proof}
$\Omega$ is measurable, being open.  Write $K=\operatorname{supp}\rho$, so
$K=\overline{B_d(0)}$ by (H3), and note that $\rho$ and $d\rho$ both vanish
wherever $d<|z|$: such a $z$ lies outside $K$, where $\rho$ vanishes together
with all its derivatives, $\rho$ being $C^\infty$ and $K$ closed.

\emph{Step A: differentiating the convolution.}  By (H3) the kernel $\rho$ is
$C^\infty$ with compact support, hence so is $d\rho$, and by (H2) $\varphi$ is
locally integrable; therefore the convolution $\varphi*d\rho$ exists and
$\varphi*\rho$ is differentiable with $D(\varphi*\rho)=\varphi*(d\rho)$
(Preamble, convolution with a smooth compactly supported kernel --- the
smoothness statement there is proved by differentiating under the integral
sign, which is what is used).  Evaluating the derivative at the basis vector
$e_i$,
\[
D_i(\varphi*\rho)(x)=\int_{\R^n}\varphi(y)\,\bigl(d\rho(x-y)e_i\bigr)\,dy .
\tag{A}
\]

\emph{Step B: the integrand vanishes off $\Omega$.}  Let $y\notin\Omega$.  Then
$y\notin\overline{B_d(x)}$ by (H4), i.e.\ $d<|x-y|$, so $\rho(x-y)=0$ and
$d\rho(x-y)=0$ by the remark above.  Hence the integrand of (A) vanishes off
$\Omega$, and the integral over $\R^n$ equals the integral over $\Omega$:
\[
D_i(\varphi*\rho)(x)=\int_\Omega\varphi(y)\,\bigl(d\rho(x-y)e_i\bigr)\,dy .
\tag{B}
\]
The same argument applies to the right-hand side's integrand, which carries the
factor $\rho(x-y)$.

\emph{Step C: the weak-gradient identity.}  Put $\psi(y)=\rho(x-y)$.  Then
$\psi$ is $C^\infty$, being $\rho$ composed with the smooth map
$y\mapsto x-y$; its closed support is the reflection of $K$ translated by $x$,
that is $\overline{B_d(x)}$, which is compact and, by (H4), contained in
$\Omega$.  So $\psi$ is an admissible test function for
\noderef{HasWeakGradientOn}, and (H2) in the direction $i$ gives
\[
\int_\Omega\varphi(y)\,D_i\psi(y)\,dy=-\int_\Omega g_\varphi(y)_i\,\psi(y)\,dy .
\]
By the chain rule $D_i\psi(y)=-\,d\rho(x-y)e_i$, the minus sign coming from
differentiating $y\mapsto x-y$; substituting that on the left and cancelling the
two signs turns the display into
\[
\int_\Omega\varphi(y)\,\bigl(d\rho(x-y)e_i\bigr)\,dy
=\int_\Omega g_\varphi(y)_i\,\rho(x-y)\,dy .
\tag{C}
\]

\emph{Step D: re-extending to $\R^n$.}  On $\Omega$ the integrand of the
right-hand side of (C) is unchanged by inserting the indicator, since
$\mathbf1_\Omega(y)=1$ there; and off $\Omega$ the indicator makes it vanish.
Hence
\[
\int_\Omega g_\varphi(y)_i\,\rho(x-y)\,dy
=\int_{\R^n}\mathbf1_\Omega(y)\,g_\varphi(y)_i\,\rho(x-y)\,dy
=\bigl(\mathbf1_\Omega(D_i\varphi)\bigr)*\rho\,(x),
\]
the first equality by splitting the integral over $\Omega$ and its complement
and the second by the definition of the convolution.  Chaining this with (B) and
(C) is the assertion.
\end{proof}
